% !Mode:: "TeX:UTF-8"

% 中英文摘要
\begin{cabstract}
本文主要系统介绍 Calabi‑Yau 簇形变理论的两个主要途径：复几何方法与代数方法。在复几何方面，我们从 Maurer–Cartan 方程出发，建立了 Kodaira–Spencer 映射，并详细讨论了 Tian–Todorov 引理在求解该方程时的关键作用，由此导出光滑 Calabi‑Yau 流形的无障碍定理。在代数方面，我们构造了 Lichtenbaum–Schlessinger 的 \(T^i\) 函子与余切复形，给出了概形形变的障碍理论，证明了形变函子的 Schlessinger 判别法及 pro‑可表示性，并在此基础上介绍了 Ran–Kawamata 的 \(T^1\) 提升原理及其在证明光滑复射影 Calabi‑Yau 簇形变函子无障碍性上的应用。
\end{cabstract}

\begin{eabstract}
This paper provides a systematic introduction to the deformation theory of Calabi--Yau varieties via two main approaches: complex geometric methods and algebraic methods. 
In the complex geometric aspect, starting from the Maurer--Cartan equation, we establish the Kodaira--Spencer map and discuss in detail the crucial role of the Tian--Todorov lemma in solving this equation, thereby deriving the unobstructedness theorem for smooth Calabi--Yau manifolds. 
In the algebraic aspect, we construct Lichtenbaum--Schlessinger's \(T^i\) functors and the cotangent complex, present the obstruction theory for deformations of schemes, prove Schlessinger's criterion and pro‑representability of the deformation functor, and on this basis introduce the Ran--Kawamata \(T^1\) lifting property together with its application to the proof of unobstructedness of deformation functor of smooth complex projective Calabi--Yau varieties.
\end{eabstract}